% ============================================
% Response to Reviewer 2
% ============================================
\section*{Response to Reviewer 2}

We thank the reviewer for their thoughtful and constructive comments, which have significantly improved the conceptual clarity of the manuscript. Below we address each point.

% ============================================
\subsection*{R2-1. Nutrient enrichment versus interaction strength}

\reviewercomment{The central claim of the paper is that increasing nutrient supply strengthens interactions, which in turn drives a shift toward community-level selection. While this is an appealing and intuitive interpretation, it may be overly simplified.

Nutrient enrichment simultaneously affects multiple ecological processes, including:
\begin{itemize}
\item species' carrying capacities
\item metabolic rates
\item environmental modification (e.g.\ pH shifts)
\item the balance of competition and facilitation
\end{itemize}

Indeed, the manuscript itself shows that pH modification by dominant taxa may contribute to winner identity in the nutrient-rich regime, while not serving as a sole explanation for Dominance frequency. This suggests that the observed transition may reflect changes in environmental feedbacks and dominance structure, rather than a simple increase in pairwise interaction strength.

This raises an important conceptual point:

\emph{Is the observed shift in coalescence outcomes driven by increased interaction strength per se, or by other effects of resource enrichment?}

We therefore suggest reframing this result more broadly in terms of changes in \emph{interaction intensity and environmental feedbacks}, rather than interpreting nutrient enrichment as a direct proxy for stronger interactions. This would align well with recent work highlighting the role of environmentally mediated interactions in microbial communities [Goldford et al., 2018, Estrela et al., 2021].}

\responselabel
We agree with the reviewer that the nutrient-gradient result should not be interpreted as nutrient enrichment directly scaling all pairwise interaction strengths. Our interpretation is that nutrient enrichment perturbs net interaction intensity and environmental feedbacks, which are reflected experimentally in increased invasion resistance and in the shift in coalescence outcomes. Multiple mechanisms, including resource competition, increased metabolic activity, pH-mediated environmental feedbacks, and changes in the competition--facilitation balance, can contribute to this nutrient-dependent change.

To make this distinction clear, we now separate the gLV interaction-strength parameter from the experimental nutrient manipulation. The gLV parameter is framed as a phenomenological parameter for ecological coupling within the model, whereas the nutrient manipulation in Fig.~4 is framed as an empirical perturbation of net interaction intensity and environmental feedbacks. We therefore use failed pairwise invasion frequency as an operational measure of invasion resistance, rather than as a direct measurement of one biochemical mechanism or a literal estimate of every pairwise Lotka--Volterra coefficient. This framing is consistent with prior work on nutrient-mediated microbial interactions and environmentally mediated community assembly (Ratzke et al., 2020; Goldford et al., 2018; Estrela et al., 2021; Duan et al., 2025).

\reviewercomment{From a consumer--resource perspective, it is also not mathematically self-evident that increasing resource supply monotonically increases effective pairwise interaction strengths. In our recent theoretical work [Duan et al., 2025], we show that when mapping mechanistic consumer--resource dynamics to effective Lotka--Volterra coefficients, the external supply rate can cancel out under broad conditions, provided that resources are not strictly limiting and metabolic strategies remain fixed.

Under such conditions, resource enrichment increases total biomass but does not necessarily increase mean per-capita competition strength. Instead, enrichment restructures the interaction network through changes in equilibrium abundances and environmental feedbacks. This distinction is important, as it suggests that the observed transition may reflect a reorganisation of interaction structure rather than a simple scaling of interaction strength.

A brief discussion of these issues would significantly strengthen the conceptual clarity of the manuscript.}

\responselabel
The apparent discrepancy comes from how the interaction coefficient is defined. The consumer--resource result discussed by Duan, Pawar and colleagues concerns coefficients derived mechanistically from resource dynamics as effective \emph{per-capita} terms in a population-growth equation. We therefore do not interpret nutrient enrichment as a monotonic increase in those effective \emph{per-capita} Lotka--Volterra coefficients. In contrast, our gLV coefficients are used as coarse-grained effective parameters in a normalized-abundance model, and our nutrient-gradient experiment uses a \emph{per-species} invasion-resistance readout. As the reviewer correctly notes, resource enrichment can increase total biomass without simply scaling the underlying \emph{per-capita} coefficients. Thus, the nutrient-gradient result should be interpreted as an increase in empirical invasion resistance, not as a claim that nutrient supply directly or monotonically scales every effective \emph{per-capita} interaction coefficient.

Results \S2.4 (``Nutrient-dependent interaction intensity and environmental feedbacks recapitulate model predictions'') now reads: \mschange{``Following prior work showing that nutrient concentration can intensify microbial competition and environmentally mediated inhibition (Hu et~al.\ 2022; Ratzke et~al.\ 2020; Hu et~al.\ 2025), we conducted additional coalescence experiments by removing or augmenting glucose and urea in the Base medium used in Fig.~1 (Methods). We treat this manipulation as a perturbation of net interaction intensity and environmental feedbacks, not as a direct measurement of a single pairwise interaction coefficient.''}

We also added a Discussion caveat addressing the consumer--resource concern directly: \mschange{``Our experimental design uses nutrient concentration to perturb net interaction intensity and environmental feedbacks, but this should not be read as a direct monotonic mapping from external resource supply to effective pairwise Lotka--Volterra coefficients. The gLV parameter $\mu$ is a simplified scalar model parameter, and the mapping between nutrient concentration and $\mu$ is necessarily approximate. This phenomenological mapping is intended to capture the net interaction intensity expressed in invasion outcomes, rather than to identify a unique biochemical mechanism. In our experiments, the nutrient-dependent increase in invasion resistance could reflect several non-mutually exclusive processes, including denser resource competition, increased metabolic activity, environmental modification such as pH shifts, changes in carrying capacities, and changes in competition--facilitation balance.''}

% ============================================
\subsection*{R2-2. Alternative explanations for community-level selection}

\reviewercomment{The manuscript interprets dominance of one parental community and correlated persistence of its taxa as evidence for community-level selection. We find this interpretation interesting and potentially insightful, but it is not uniquely diagnostic on its own.

Similar patterns could arise from:
\begin{itemize}
\item shared environmental filtering
\item correlated traits within parental communities
\item environmental modification (e.g.\ pH tolerance)
\end{itemize}

In particular, the nutrient-rich regime appears consistent with strong environmental filtering imposed by a small number of dominant taxa, which may not require invoking selection acting at the level of the community as a unit.

We therefore recommend clarifying the interpretation of ``community-level selection,'' and explicitly discussing alternative mechanisms that could generate similar patterns. This distinction has been discussed in previous work on community coalescence and microbial assembly [Mansour et al., 2018, Rillig et al., 2015].}

\responselabel
We agree that Dominance and correlated persistence are not uniquely diagnostic of one mechanism. We have therefore revised the manuscript to present Dominance as phenomenological support for community-level selection, rather than as direct evidence for one causal process. We agree that multiple mechanisms, including shared environmental tolerance, correlated traits, and filtering by pH-modifying dominant taxa, can generate this outcome-level pattern; Dominance alone does not prove one exclusive mechanism.

We made this distinction explicit in Results \S2.1 (``Community-level selection is prevalent in microbial coalescence''): \mschange{``Crucially, Dominance provides phenomenological support for community-level selection: one parental community displaces the other while retaining its internal compositional structure, indicating that its species persist collectively.''}

This clarification is particularly important for Nutr$+$. We agree that Nutr$+$ is consistent with a top-down/environmental-feedback route in which dominant taxa may bias persistence toward their parental community, but pH contrast alone does not explain the full Dominance pattern: acid--alk pairings were not significantly enriched for Dominance relative to pairings in which both parental communities had similar pH in either Base or Nutr$+$ (Fisher's exact tests: Base $p = 0.49$, Nutr$+$ $p = 0.47$; Response Fig.~R1-2). In acid--alk coalescence pairs, the acidic parental community wins $86.4\%$ of Nutr$+$ events (38/44, binomial $p = 9.4 \times 10^{-7}$; Response Fig.~R1-2). By contrast, Base medium is framed as the stronger case for an emergent multi-species route, because dominant-species predictability is much weaker there.

We revised the Discussion paragraph defining community-level selection accordingly. It now states, in part: \mschange{``We define community-level selection operationally as origin-correlated persistence during coalescence, rather than as a claim about one exclusive biochemical mechanism. This is a phenomenological description of the outcome, not a mechanistic attribution on its own. Under this definition, shared environmental tolerances, correlated traits, and environmental filtering by dominant pH-modifying taxa are alternative mechanistic routes that can produce the same outcome-level pattern when they make parental-community origin predictive of persistence. Our scope here is to define and test community-level selection as an outcome-level pattern, while future work can directly disentangle how these mechanisms contribute across different coalescence regimes.''}

To test environmental feedback explicitly, we ran a Ratzke--Gore-style pH-feedback model (Ratzke and Gore, 2018), which produces asymmetric replacement over an intermediate range of feedback strength (Response Fig.~R2-2). We agree that pH-mediated environmental modification is a plausible mechanism for coalescence outcomes, and we do not claim that the effective-interaction gLV is the unique explanatory framework. One observation from this comparison is that the pH-only model rarely produces Restructuring outcomes ($\sim 1\%$ across feedback strengths; 3/297 simulated events), whereas Restructuring is observed at substantially higher rates both in our experiments ($\sim 8\text{--}31\%$ across nutrient conditions) and in the effective-interaction gLV ($\sim 22\%$ across $\mu$ values). As discussed in R2-1, our gLV interaction parameter is treated as an effective coupling that can incorporate environmental-feedback contributions including pH modification; under this reading, pH-feedback and effective-interaction-strength descriptions are complementary rather than mutually exclusive.

\begin{figure}[H]
\centering
\includegraphics[width=0.88\textwidth]{internal_Q5_phase_pH.pdf}
\caption*{\textbf{Response Fig.~R2-2.} \textbf{Outcome distribution of a Ratzke--Gore-style pH-feedback model across feedback strengths.} Per-event outcomes in the similarity plane (left, three feedback strengths) and stacked outcome fractions (right) for the pH-feedback model. The model produces asymmetric replacement (Dominance) at moderate-to-strong feedback, supporting pH-mediated environmental modification as a plausible mechanism contributing to coalescence outcomes. Restructuring outcomes are rare across all feedback strengths ($\sim 1\%$; 3/297 simulated events), in contrast to the substantial Restructuring fractions observed experimentally ($\sim 8\text{--}31\%$) and in the effective-interaction gLV ($\sim 22\%$).}
\end{figure}

In addition, several related comments raise alternative theoretical models and complementary controls, including biomass effects, pH-mediated filtering, dominant-species circularity, pool size, richness, and facilitation, which are addressed in R1-1 to R1-4 and R3-2 to R3-4 and consolidated in Supplementary Note~5 (``Alternative models and controls for community-level selection''). The passive-retention alternatives are most relevant here: the denser parental community does not generally win, and we did not detect a significant overall effect of initial pool size on Dominance frequency (R1-1, R1-4).

\reviewercomment{Relatedly, we suggest clarifying the ecological meaning of ``failed invasion.'' At present it is primarily used as a proxy for interaction strength, but it is also naturally interpretable as invasion resistance (i.e.\ whether a rare invader can establish). Framing this more explicitly in terms of invasion fitness would help connect the experimental results to ecological theory and may provide additional insight into the observed patterns.}

We thank the reviewer for this important clarification. We agree that failed invasion is naturally connected to invasion resistance and invasion fitness. We therefore revised the manuscript to describe the pairwise assay as an operational measure of invasion resistance and net interaction intensity across nutrient conditions. In Fig.~4B, we retained the y-axis as the measured quantity, ``frequency of failed invasions,'' while the panel framing identifies this as an invasion-resistance analysis. Results \S2.4 now states: \mschange{``To operationally quantify its effect, we measured failed invasion frequency using pairwise invasion assays among the 12 most abundant isolates (95:5 initial frequency; Methods, Supplementary Figs.~5--7). The fraction of failed invasions increased monotonically with nutrient supply (Nutr$-$: $2 \pm 1\%$; Base: $33 \pm 4\%$; Nutr$+$: $48 \pm 4\%$; mean $\pm$ s.e.m.; Fig.~4B), indicating that invasion resistance increases across the nutrient gradient.''}

The Methods now states that \mschange{``the fraction of failed invasions served as an empirical proxy for net interaction intensity and invasion resistance.''} We distinguish this two-species assay from the pairwise selection-correlation metric in R2-4.

% ============================================
\subsection*{R2-3. Continuous measures alongside categorical outcomes}
\reviewercomment{The classification into Dominance, Mixture, and Restructuring is useful and provides a clear framework for summarising outcomes. However, it depends on thresholding a continuous similarity measure.

As noted in the broader literature [Mansour et al., 2018], coalescence outcomes often lie along a continuum rather than discrete categories.

We suggest that the authors:
\begin{itemize}
\item present continuous similarity measures alongside categorical outcomes
\end{itemize}}

\responselabel
We agree that the Dominance / Mixture / Restructuring classifications depend on boundaries in a continuous similarity space. PDI was already introduced in Figs.~3 and 4 as a continuous measure of how strongly the coalesced community is directionally weighted toward one parental community versus the other. To make the thresholding explicit, we now present PDI together with asymmetry and retention magnitude as continuous auxiliary metrics underlying the discrete outcome classification, as shown in Response Fig.~R2-3, and have added these analyses to the Supplementary Information as Supplementary Figs.~29--31. These measures show where each coalescence event lies before category boundaries are applied, while the categorical labels provide the compact summary used in the main analyses.

To verify how sensitive our main nutrient-dependent result is to the classification-boundary choices, we varied the PDI and retention boundaries around the manuscript baseline (Response Fig.~R2-4). Across a broad range of boundary choices, the Dominance fraction remains ordered Nutr$-$ $<$ Base $<$ Nutr$+$, supporting the robustness of the nutrient-dependent trend. We also added this boundary-sensitivity analysis to the Supplementary Information as Supplementary Fig.~32 and refer to it in Supplementary Note~1.

\begin{figure}[H]
\centering
\includegraphics[width=0.9\textwidth]{marginal_distributions_by_medium.pdf}
\caption*{\textbf{Response Fig.~R2-3.} \textbf{Continuous similarity measures track the same nutrient-dependent shift as the categorical classifications.} Rows show Nutr$-$, Base, and Nutr$+$; columns show Parental Dominance Index (PDI), where values near 0 or 1 indicate strong asymmetry toward one parental community and values near 0.5 indicate balanced parental contributions; asymmetry magnitude $y$, a direction-independent measure of how unequal the two parental contributions are; and squared retention magnitude $x^2 = u^2 + v^2$, which measures how much of the coalesced composition is explained by the parental communities.}
\end{figure}

\begin{figure}[H]
\centering
\includegraphics[width=0.95\textwidth]{boundary_sensitivity_by_medium.pdf}
\caption*{\textbf{Response Fig.~R2-4.} \textbf{Boundary-sensitivity analysis supports the nutrient-dependent Dominance trend.} (A) Dominance fractions after varying the PDI boundary while holding the retention boundary at $x^2 > 0.5$. (B) Dominance fractions after varying the retention boundary while holding the PDI boundary at the manuscript baseline. Dashed lines mark the manuscript baseline. This sensitivity analysis is descriptive.}
\end{figure}

\reviewercomment{In addition, the distinction between mixture and restructuring is not entirely straightforward biologically, particularly because abundance changes contribute to classification.

We suggest that the authors:
\begin{itemize}
\item clarify the biological interpretation of ``restructuring''
\end{itemize}}

We clarify that Restructuring denotes low parental retention: the coalesced community is not well explained by either parental community alone or by a non-negative linear combination of the two parental communities. In this interpretation, Restructuring can arise when cross-community species pairings that were not present during parental assembly generate a new stable composition, especially in complex multi-species communities where abundance shifts and higher-order interactions make the final state poorly described by the parental-community basis.

Results \S2.1 (``Community-level selection is prevalent in microbial coalescence'') now includes the concise Restructuring definition: \mschange{``Restructuring denotes low parental retention: the coalesced community is not well explained by either parental community alone or by a balanced linear combination of the two parental communities, and may reflect a new stable composition produced by previously unseen cross-community species interactions.''}

\reviewercomment{We also find the comparison between dot product and Jaccard metrics in Extended Data Fig.~2 potentially very informative. The divergence between these metrics in low-nutrient environments may itself provide evidence that species-level processes dominate in that regime, and this could be discussed more explicitly.}

We agree that this comparison is informative. Extended Data Fig.~2 was generated for the Base-medium outcomes, so we discuss the metric divergence in that context. Consistent with our response to Reviewer~1, we now frame the result as a difference in what each metric measures rather than as a broad robustness claim. This pattern is consistent with the fact that the metrics emphasize different aspects of community composition: relative-abundance metrics are more sensitive to quantitative dominance in relative abundance, whereas Jaccard reflects presence/absence retention and Jensen--Shannon evaluates divergence across the full abundance distribution. Thus, a coalesced community can be skewed toward one parental community in abundance space while still retaining taxa from both parental communities, or while redistributing subdominant taxa in ways that make it less skewed toward one parental community under Jaccard or Jensen--Shannon.

Results \S2.1 now narrows the main-text robustness statement to the two alternative metrics that recovered the same pattern: \mschange{``This pattern of Dominance as the most frequent outcome was robust under similarity metric choices including Euclidean distance and Bray--Curtis dissimilarity (Extended Data Fig.~2).''} Extended Data Fig.~2 caption now carries the metric-divergence interpretation: \mschange{``The primary Dominance pattern is recovered under similarity metric choices including Euclidean distance and Bray--Curtis dissimilarity, but not under Jensen--Shannon divergence or Jaccard index, for which the coalesced community is less often measured as skewed toward one parental community. This divergence reflects the fact that the metrics emphasize different aspects of community composition: vector decomposition, Euclidean distance, and Bray--Curtis dissimilarity capture quantitative similarity to each parental community, whereas Jaccard reflects presence/absence retention and Jensen--Shannon evaluates divergence across the full abundance distribution.''}

% ============================================
\subsection*{R2-4. Interpretation of pairwise selection correlation}

\reviewercomment{The pairwise ``selection correlation'' metric is an interesting attempt to quantify lineage-level coherence, but its interpretation remains somewhat unclear.

It appears closely related to co-occurrence patterns, and it is not fully clear whether the observed correlations reflect ecological interactions, shared environmental responses, or methodological effects. We find this analysis potentially insightful, but its conceptual interpretation would benefit from clarification.}

\responselabel
We agree that the pairwise selection correlation needed a clearer interpretation. Operationally, it tests, for each coalescence event, whether species pairs from the same parental community share survival outcomes more strongly than species pairs from different parental communities. The comparison is evaluated against a within-event null in which parental-origin labels are shuffled. Generic co-occurrence, any tendency for species to share survival outcomes under a given environment or coalescence condition, is absorbed by the within-event null and therefore does not contribute to the same-parental-community enrichment. We use this enrichment as evidence for community-level selection, while recognizing that the metric does not identify the causal route; direct competition, shared environmental tolerance, and environmental modification can all generate such structure.

Results \S2.2 now defines the metric directly: \mschange{``We quantified this effect using pairwise selection correlation, which tests across coalescence events whether species pairs from the same parental community share survival outcomes (both persist or both go extinct) more strongly than species pairs from different parental communities (Supplementary Information; Fig.~2D).''}

\reviewercomment{We suggest adding a short paragraph linking this metric more explicitly to invasion fitness in a gLV framework. In that context, invasion fitness corresponds to the growth rate of a species when rare in a resident community. The pairwise invasion assays can then be interpreted as empirical approximations of invasion fitness in two-species systems, while coalescence outcomes reflect correlations in invasion success across many species.

This would provide a unifying theoretical framework linking pairwise assays, community coalescence, and the notion of community-level selection.}

\responselabel
We agree that this point needs clarification, especially because ``pairwise'' refers to two different analyses in the manuscript. As discussed in R2-2, the 95:5 pairwise invasion assay measures growth dynamics and ecological outcomes in two-species cocultures, and we use it primarily as an empirical approximation of invasion resistance and net interaction strength under each nutrient condition. By contrast, the pairwise selection-correlation metric is computed from community coalescence outcomes. It quantifies whether assembly history creates correlated species fates across species pairs: specifically, whether species pairs from the same parental community show more strongly coupled survival or extinction than species pairs drawn from different parental communities.

Finally, we acknowledge the limitations and scope of this interpretation more explicitly in the Discussion: \mschange{``We define community-level selection operationally as origin-correlated persistence during coalescence, rather than as a claim about one exclusive biochemical mechanism.''} \mschange{``Accordingly, our evidence for community-level selection is operational and outcome-based: we assess it from parental-community similarity after coalescence and pairwise species-fate correlation, rather than directly measuring the mechanistic selection pressures that generate these outcomes.''} \mschange{``Thus, our scope here is to define and test community-level selection as an outcome-level pattern, while future work can directly disentangle how different mechanisms contribute across coalescence regimes.''}

% ============================================
\subsection*{R2-5. Pre-selection of natural communities}

\reviewercomment{The inclusion of natural communities is an important strength of the study. However, the use of a common stabilisation phase in defined media may introduce pre-selection effects that reduce ecological heterogeneity across communities.

In particular, incubating diverse natural communities in simplified media is known to drive convergence toward a limited set of functional guilds [Goldford et al., 2018]. As a result, the ``natural'' communities used here may already be filtered to resemble the synthetic communities, potentially contributing to the similarity of results across systems.

We therefore suggest that the authors:
\begin{itemize}
\item discuss the potential effects of this pre-selection more explicitly
\item clarify the extent to which taxonomic or functional convergence occurs during the stabilisation phase
\end{itemize}

This would help contextualise the generality of the natural-community results.}

\responselabel
We agree this is an important caveat. As the reviewer notes, defined-media enrichment can drive taxonomic or functional convergence (Goldford et al., 2018), and we cannot directly test this here because the present experiments collected 16S community profiles rather than metagenomic, metabolomic, or trait-resolved measurements. We therefore narrow our claim: the natural sample-derived experiment shows that the nutrient-dependent increase in Dominance also occurs in taxonomically richer, laboratory-stabilized communities derived from natural samples, not that unmanipulated natural ecosystems follow the same rules.

Within the available data, we asked whether stabilized communities still carried source-specific taxonomic signal. Using same-source biological replicates as the empirical retention baseline, different-source parental communities remained substantially less similar at the ASV level (Jaccard similarity 1.7- to 4.1-fold lower than same-source replicates across the three media; Response Fig.~R2-5b; Supplementary Fig.~44b), with a comparable pattern at the Genus level (Response Fig.~R2-5e; Supplementary Fig.~44e). We acknowledge that distinguishability between sources does not rule out the Goldford-style scenario in which each source converges to its own source-specific attractor during stabilization; only a pre-enrichment baseline could test that directly, and we do not have one. The data do, however, exclude the stronger version of pre-selection in which all sources collapse to a single shared taxonomic endpoint. Consistent with this, coalesced communities also did not collapse to a shared within-medium endpoint, and each remained closer to its own parental-community pair than to unrelated same-medium parental communities (own-parent ASV Jaccard similarity exceeded unrelated-parent similarity in all three media, $n=30$ events per medium, one-sided paired Wilcoxon signed-rank $p<10^{-6}$; Response Fig.~R2-5b,c,e,f; Supplementary Fig.~44b,c,e,f).

We therefore maintain the claim that nutrient-dependent Dominance occurs in laboratory-stabilized communities of higher initial taxonomic diversity derived from natural samples, while noting that source-specific or functional pre-selection during the seven-cycle enrichment may attenuate the extent to which it generalises to unmanipulated natural ecosystems. This revision also addresses the related Reviewer~3 concern that natural sample-derived communities are still model systems and require cautious claims (R3-4).

\begin{figure}[H]
\centering
\includegraphics[width=0.95\textwidth]{Fig_R2_5_natural_taxonomic_distinctness.pdf}
\caption*{\textbf{Response Fig.~R2-5.} \textbf{Natural sample-derived communities do not collapse to a single medium-specific taxonomic endpoint after stabilization or coalescence.} \textbf{a,d}, ASV- and Genus-level richness of natural parental and coalesced communities by medium. \textbf{b,e}, Pairwise Jaccard similarity among same-source stabilized parental replicates (S; empirical source-retention baseline; $n=6$ pairs per medium), different-source stabilized parental communities (D; post-stabilization convergence check; $n=60$ pairs per medium), and final coalesced communities (C; outcome-level distinctness check; $n=435$ pairs per medium) within each medium, computed at the ASV and Genus levels; per-box sample sizes are annotated above each box. \textbf{c,f}, For each final coalesced community ($n=30$ coalescence events per medium), mean similarity to its own two parental communities versus unrelated stabilized parental communities from the same medium, measured by ASV- and Genus-level Jaccard similarity; own-parent similarity exceeded unrelated-parent similarity in every medium (one-sided paired Wilcoxon signed-rank test, all $p<10^{-6}$ at both the ASV and Genus levels), with per-medium significance ($p<0.001$) and event counts annotated in each panel. Analyses use the raw natural ASV table, matched taxonomy annotations, ASVs above the 0.1\% relative-abundance threshold, and post-stabilization 16S profiles; unresolved higher-rank taxonomy labels are kept ASV-specific when ASVs are collapsed by rank. These analyses test for complete post-stabilization taxonomic collapse and retained parent-specific taxonomic signal, but not pre-to-post enrichment convergence or functional convergence.}
\end{figure}

Manuscript changes: we narrowed the natural-community framing throughout the manuscript, including the Abstract, Introduction, Methods, Results subsection title, and Fig.~6 caption, so the result is described as applying to natural sample-derived communities established through laboratory stabilization rather than unfiltered natural ecosystems. We also added the same-source/different-source taxonomic-distinctness analysis to the Supplementary Information as Supplementary Fig.~44, and added a separate paragraph to Results \S2.6 to present pre-selection during stabilization as an alternative explanation for similarity between synthetic and natural sample-derived outcomes:
\mschange{``These experiments tested whether nutrient-dependent Dominance is also observed after laboratory assembly of taxonomically richer communities from natural samples. An alternative explanation is that the seven-cycle enrichment phase imposed laboratory pre-selection before coalescence. Under this hypothesis, communities that began from natural samples could have converged toward taxa or functional strategies favored by the defined media, making the subsequent coalescence outcomes resemble those of synthetic communities. The higher post-stabilization richness and low ASV overlap among communities from different source samples argue against complete ASV-level convergence (Supplementary Fig.~44), but the current data cannot determine how much of the similarity between synthetic and natural sample-derived outcomes reflects intrinsic natural-community dynamics versus pre-selection during stabilization.''}
We also revised the Discussion caveat to state explicitly that pre-selection may occur and that the current data cannot quantify taxonomic convergence during stabilization or functional convergence.


% ============================================
\subsection*{R2-6. Frame gLV as phenomenological}

\reviewercomment{The gLV model successfully reproduces the qualitative transition in coalescence outcomes and serves as a useful minimal framework. However, it is limited to purely competitive interactions and does not explicitly capture environmentally mediated effects such as pH modification.

We therefore suggest framing the model more explicitly as a phenomenological rather than mechanistic description of the system.}

\responselabel
We agree that the gLV model should be framed as a phenomenological competitive baseline rather than as a mechanistic description of the biochemical interactions operating in the experiments. Results \S2.2 (``Theoretical model with random interactions reproduces community-level selection'') now states: \mschange{``The gLV model serves as a phenomenological framework to explore how the statistical properties of interaction coefficients shape coalescence outcomes, rather than as a mechanistic model of the specific biochemical interactions (e.g., pH modification, metabolic cross-feeding) underlying our experimental observations.''}

\reviewercomment{In addition, while the robustness analysis across interaction coefficient distributions is appreciated, a brief discussion of the biological interpretation of these distributions would be helpful. In particular, clarifying how the mean interaction strength parameter ($\mu$) relates to underlying ecological processes would improve interpretability.}

\responselabel
Biologically, $\mu$ parameterizes the sampling range for competitive interaction coefficients in the gLV model. Because our baseline modeling framework draws $\alpha_{ij} \sim U(0, 2\mu)$, increasing $\mu$ raises both the mean inhibitory effect of $\alpha_{ij}$ and the variance of sampled coefficients. The first effect pulls the system away from near-neutral dynamics and produces appreciable growth inhibition. The second creates stronger mismatches between species that assembled together and outsiders, which lets assembly filter species into mutually compatible groups, the ``cohesion without cooperation'' route to Dominance discussed in R3-4.

We decoupled these two effects in Supplementary Note~3 and report the empirical mean-vs-variance result in R3-4 (Response Fig.~R3-4f). Both coefficient mean and coefficient variance contribute to Dominance (Pearson $r=0.53$ and $0.78$ respectively), and neither alone is sufficient. This is consistent with the Hu et al.\ framework, in which $\mu$ couples mean and variance simultaneously through the uniform distribution $U(0, 2\mu)$.

We also revised the Results, Methods, and Supplementary Methods to avoid presenting $\mu$ as simply a mean interaction strength and to define it instead as the sampling-range parameter for competitive interaction coefficients. Results \S2.2 now states: \mschange{``We fixed $r_i = 1$ and drew off-diagonal competition coefficients from a uniform distribution $\mathbb{U}(0, 2\mu)$, so $\mu$ controls the width of the interaction-coefficient sampling distribution and therefore increases both coefficient mean and variance: higher $\mu$ means stronger and more heterogeneous interspecies competition.''} Methods now states: \mschange{``We fixed growth rates to 1 and self-interaction coefficients $\alpha_{ii} = 1$, and drew off-diagonal competition coefficients from a uniform distribution $\mathbb{U}(0, 2\mu)$, so $\mu$ sets the range of sampled competitive interaction coefficients and thereby controls both their mean and variance within this phenomenological model.''} Supplementary Methods further defines $\alpha_{ij}$ as an effective per-capita inhibitory term and states that the baseline $\alpha_{ij}\geq 0$ model captures competitive growth inhibition rather than facilitation.


% ============================================
\subsection*{R2 minor comments}

\reviewercomment{Improve the visualisation of pairwise correlation results, as the separation between groups is not visually clear.}

\responselabel
We agree that the previous visualization could be clearer. We revised the Fig.~2D caption in Results \S2.2 to define each visual element explicitly:
\mschange{``Individual dots show per-event pairwise selection correlations; in the simulation panel, 100 stratified events per group are displayed for legibility. Squares with error bars show mean $\pm$ s.e.m.\ across all coalescence events ($n = 1{,}200$ per group in simulations and all valid experimental events in the experimental panel). Gray horizontal lines indicate the random selection baseline from a null model in which species origin labels are shuffled (see Supplementary Information).''}

\reviewercomment{Specify how pH was measured (e.g.\ continuously or at endpoints), and whether variability across replicates was assessed.}

We clarified the pH-measurement description in Supplementary Methods, ``Optical Density Measurements and Monoculture Characterization,'' which now specifies the timing and replicate treatment:
\mschange{``Community-level pH measurements were endpoint measurements on culture supernatant rather than continuous measurements during growth; replicate-to-replicate variability is represented by the biological replicates shown in the pH-based analyses.''}

\reviewercomment{Clarify the biological interpretation of the interaction coefficient distributions used in the theoretical model.}

We clarified the biological interpretation of the gLV coefficients and their distributions in Supplementary Methods, ``Lotka--Volterra Simulations.'' The revised text states:
\mschange{``Biologically, each $\alpha_{ij}$ is an effective per-capita term for the effect of species $j$ on species $i$ within the coarse-grained gLV model. It can absorb the net inhibitory consequences of direct mechanisms (e.g., resource competition, interference) and indirect mechanisms (e.g., metabolite-mediated inhibition, pH modification) at steady state, but it does not identify the underlying biochemical mechanism and the model does not dynamically represent pH or other environmental state variables.''}
We also clarify the meaning of $\mu$ and the distributional assumption:
\mschange{``Thus, in the baseline competitive ensemble, $\mu$ sets the range of sampled competitive interaction coefficients and thereby controls both their mean and variance. Increasing $\mu$ simultaneously increases the average competitive inhibition and the heterogeneity of pairwise coefficients, rather than isolating either effect.''}
\mschange{``This distribution is a phenomenological ensemble of heterogeneous effective competitive effects rather than a fitted empirical distribution of biochemical mechanisms; robustness analyses across alternative coefficient distributions and a separate mean-vs-variance sweep identify which aspects of interaction heterogeneity are required for the qualitative coalescence transition.''}

\reviewercomment{Ensure consistent terminology and notation throughout (e.g.\ ``pairwise'', ``community-level'', $\alpha_{ij}$).}

We standardized the terminology throughout the revised manuscript. In particular, Results \S2.2 now uses ``pairwise selection correlation'' for the lineage-level correlation metric and defines it as follows:
\mschange{``We quantified this effect using pairwise selection correlation, the degree to which species pairs share the same survival outcome (both persist or both go extinct) across coalescence events.''}
We use ``community-level selection'' for origin-correlated persistence at the parental-community level and use $\alpha_{ij}$ consistently for the gLV interaction coefficients.

\reviewercomment{Correct minor typographical issues, including ``generalis ability'' $\to$ ``generalisability''.}

We corrected ``generalis ability'' to ``generalisability'' and addressed other minor typographical issues throughout the revised manuscript.
